\section{The transition on the large-root side}\label{sec:bridge}

Put $X=e^u>1$ and
\[
 W_u(k,h)=k^n h^{n-k+1}X^{-(k+h-2)},\qquad k,h\ge1.
\]
The positive double sum converges, since
$W_u(k,h)\le k^nh^nX^{-(k+h-2)}$. Eulerian symmetry gives
\begin{equation}\label{eq:bridge-exact}
 \frac{P_n(-X)}{(-X)^n}
 =\sum_{\substack{k,h\ge1\\k+h\le n+1}}
 (-1)^{k+h-2}A(n,k)A(n-k+1,h)X^{-(k+h-2)}.
\end{equation}
Lemma~\ref{lem:bonferroni} bounds every absolute summand by $W_u(k,h)$.
Throughout this section, $\rho_j=(j+1)/j$.

\subsection{The path of maximizing cells}

For an integer $m\ge2$, define
\[
 f(k,h)=n\log k+(n-k+1)\log h,\qquad
 M_m=\max_{k+h=m}e^{f(k,h)}.
\]
Choose the smallest maximizing $k=k_m$, and set $h_m=m-k_m$.
For $2\le m<m+1\le\lfloor\sqrt n\rfloor$, this path increases exactly
one coordinate by one at each step. Indeed, for $h=m-k\ge2$ its
adjacent difference is
\[
 q_m(k)=n\log\rho_k-\log h-(n-k)\log\rho_{h-1}.
\]
The comparisons
\begin{align*}
 q_{m+1}(k)-q_m(k)
   &=(n-k)\log\frac{h^2}{h^2-1}-\log\rho_h>0,\\
 q_{m+1}(k+1)-q_m(k)
   &=-n\log\frac{(k+1)^2}{k(k+2)}+\log\rho_{h-1}<0
\end{align*}
follow from $n-k>h$ and $n/(k+1)^2\ge1>\log2$, respectively.
They also imply that $q_m(k)$ strictly decreases in $k$. Its zero
crossing therefore moves by zero or one position when $m$ increases.
The smallest-maximizer convention includes a tie of adjacent cells;
boundary cases use the same one-sided comparisons.

The sequence $\log M_m$ is strictly concave. Two successive steps in
one coordinate have decreasing increments by coordinate concavity.
For a $k$ step followed by an $h$ step, optimality gives
$f(K+1,H)\ge f(K,H+1)$, whereas
\[
 [f(K,H+1)-f(K,H)]
 -[f(K+1,H+1)-f(K+1,H)]=\log\rho_H>0.
\]
Thus the second increment is strictly smaller; the opposite order is
identical. Consequently the global maximizing total degree changes
one at a time as $u$ decreases, at the ordered times
$\log(M_{m+1}/M_m)$.

The two coordinate events starting from $(K,H)$ have balance times
\begin{equation}\label{eq:bridge-balances}
 \alpha=n\log\rho_K-\log H,\qquad
 \beta=(n-K+1)\log\rho_H.
\end{equation}
Consecutive events in the same coordinate are separated by more than
\begin{equation}\label{eq:bridge-spacing}
 \frac{n}{(K+1)^2},\qquad \frac{n-K+1}{(H+1)^2},
\end{equation}
respectively. Their exact differences are
$n\log((K+1)^2/(K(K+2)))+\log(H'/H)$, or
$(n-K+1)\log((H+1)^2/(H(H+2)))+(K'-K)\log\rho_{H+1}$,
where $H'\ge H$ and $K'\ge K$ along the path.

\subsection{Global descent estimates}

Fix $(K,H)$ and put $N=n-K+1$, $S=K+H$. For integer shifts
$p\ge1-K$, $q\ge1-H$, write
\[
 w_{pq}=\frac{W_u(K+p,H+q)}{W_u(K,H)},\quad
 A=\frac{\rho_K^n}{HX},\quad B=\frac{\rho_H^N}{X},\quad
 d(p)=\max\{-p,p-1,0\}.
\]

\begin{lemma}\label{lem:bridge-joint}
Suppose $K,H\ge7$, $n-S>S^2$,
\[
 \delta_k=n\log(\rho_K/\rho_{K+1})\ge12,\qquad
 \delta_h=N\log(\rho_H/\rho_{H+1})\ge12,\qquad
 e^{-2.02}\le A,B\le e^{2.02}.
\]
With $\varepsilon=10^{-4}$, one has
\[
 w_{pq}\le57AB\varepsilon^{d(p)+d(q)},\qquad
 \sum_{(p,q)\notin\{0,1\}^2}|pq|w_{pq}<.04AB.
\]
\end{lemma}

\begin{proof}
At an outside cell decrease the coordinate with the larger distance
from $\{0,1\}$. The four outward ratios are bounded by
$Ae^{-G_+(d)}$, $A^{-1}e^{-G_-(d)}$, $Be^{-J_+(d)}$, or
$B^{-1}e^{-J_-(d)}$, where
\begin{align*}
G_+(d)&=n\log\frac{\rho_K}{\rho_{K+d}}
             -\log\frac{H}{\max(1,H-d)},&&p=d+1,\\
G_-(d)&=n\log\frac{\rho_{K-d}}{\rho_K}
             -\log\frac{H+d+1}{H},&&p=-d,\\
J_+(d)&=N\log\rho_H-[N+\min(d,K-1)]\log\rho_{H+d},&&q=d+1,\\
J_-(d)&=(N-d-1)\log\rho_{H-d}-N\log\rho_H,&&q=-d.
\end{align*}
These expressions are used for admissible cells and $d\ge1$.
For example the first case has $h\ge\max(1,H-d)$; the last has
$k\le K+d+1\le S<n$.
All four losses increase with $d$. Before truncation, their successive
increments are bounded below, respectively, by
\[
 \frac{n}{(K+d+1)^2}-\log2,\quad
 \frac{n}{(K-d)^2}-\log2,\quad
 \frac{N}{(H+d+1)^2}-\log2,\quad
 \frac{n-S}{S^2}-\log2,
\]
which are positive. For the fourth, putting $m=H-d\ge2$ gives the
exact increment
$(N-d-2)\log(\rho_{m-1}/\rho_m)-\log\rho_m$.
After truncation in $G_+$ or $J_+$ monotonicity is immediate.
The first-layer losses are at least
\[
 \delta_k-\log\frac H{H-1},\quad
 \delta_k-\log\frac{H+2}H,\quad
 \delta_h-\log\rho_{H+1},\quad
 \delta_h-2\log\rho_{H-1}.
\]
Each correction is at most $1/3$. Thus every selected outward ratio
is at most $e^{2.02-12+1/3}<\varepsilon$.
This holds on the entire positive quadrant: a selected negative
$h$ step forces $k\le S<n$, whereas in a positive $h$ step increasing
$k$ only decreases the ratio, even when $n-k+1$ is negative.

Descent takes exactly $d(p)+d(q)$ steps to the central four cells.
Their largest normalized weight is
$M=\max\{1,A,B,ABH/(H+1)\}<57AB$, since $e^{4.04}<57$.
Finally
\[
 \sum_{p\in\mathbb Z}|p|\varepsilon^{d(p)}
 =\frac{1+\varepsilon}{(1-\varepsilon)^2}.
\]
Only $(1,1)$ contributes inside the central box, so the weighted tail is
less than
\[
 57AB\left\{\left[\frac{1+\varepsilon}{(1-\varepsilon)^2}\right]^2-1\right\}
 <57AB\,7\varepsilon<.04AB.
\]
\end{proof}

\begin{lemma}\label{lem:bridge-vertex}
Suppose $K,H\ge7$, $n-K-H>(K+H)^2$, and the four coordinate-neighbor
weights are at most $e^{-2}W_u(K,H)$. Then
\[
 \sum_{(k,h)\ne(K,H)}W_u(k,h)<\frac{400}{441}W_u(K,H).
\]
If also $n>1600$ and $\max(K,H)\le\sqrt n/4+2$, the sign of
$P_n(-X)$ is $(-1)^{n+K+H-2}$ and exactly $K+H-2$ zeros lie outside
$|x|=X$, counted with multiplicity.
\end{lemma}

\begin{proof}
Use the preceding descent with distances $|p|,|q|$ from one cell.
The first-layer boosts over the corresponding coordinate-neighbor
ratios are at most $H/(H-1)$, $(H+1)/H$, $\rho_H$, $\rho_{H-1}$.
Their logarithms are at most $1/6$. Subsequent losses increase by the
same four difference calculations, now starting at the single cell.
Thus $w_{pq}\le r^{|p|+|q|}$ for $r=e^{-11/6}<4/25$. Summation gives
$\sum_{(p,q)\ne(0,0)}w_{pq}<4(4/25)/(1-4/25)^2=400/441$.
The two central Bonferroni errors sum to at most
$2(n+1)e^{-(8/5)\sqrt n}<.001$ under the stated coordinate bounds.
Every other actual weight is bounded by its ideal weight. Hence the
actual central monomial of degree $n-K-H+2$ dominates the remainder
on $|x|=X$. Rouch\'e's theorem gives the count, and evaluation at $-X$
gives the sign.
\end{proof}

\subsection{The sign inside a paired transition}

Under the hypotheses of Lemma~\ref{lem:bridge-joint}, there is an exact factorization
\[
 w_{pq}=a_pb_qc_{pq},\quad
 a_p=(1+p/K)^n(HX)^{-p},\quad
 b_q=(1+q/H)^NX^{-q},\quad c_{pq}=(1+q/H)^{-p}.
\]
Let $F(u)=\sum_p(-1)^pa_p$ and $G(u)=\sum_q(-1)^qb_q$.
All sums converge absolutely. The residual is one on both axes and
$c_{11}=H/(H+1)$. Moreover
\begin{equation}\label{eq:bridge-residual}
 |c_{pq}-1|\le\frac{|pq|}{H}\max(1,c_{pq}).
\end{equation}
Indeed, if $t=1+q/H>0$ and $m=|p|$, its left side divided by the
maximum is $1-\min(t,t^{-1})^m\le m[1-\min(t,t^{-1})]\le m|t-1|$.
Consequently
\begin{equation}\label{eq:bridge-product}
 \sum_{p,q}(-1)^{p+q}w_{pq}
 =F(u)G(u)-\frac{AB}{H+1}+E,\qquad |E|<.08AB/H.
\end{equation}
Here Lemma~\ref{lem:bridge-joint} bounds the weighted tail of $w$ by
$.04AB$. The independent product $a_pb_q$ has the same bound: its
first outward ratios are at most $e^{2.02-12}<10^{-4}$ and subsequent
ratios decrease. Sum \eqref{eq:bridge-residual} outside the central box.

We also need the actual Eulerian coefficients. Suppose $n\ge1600$
and $\max(K,H)\le\sqrt n/4+2$. In the region $d(p)<H$, $d(q)<H$,
both coordinates are at most $\sqrt n/2+4$, all cells belong to
$k+h\le n+1$, and
\[
 n/k\ge(8/5)\sqrt n,\qquad (n-k+1)/h\ge(8/5)\sqrt n.
\]
The second follows at $t=\sqrt n\ge40$ from
$t^2/5-69t/10-3>0$; the first is weaker.
Bonferroni therefore bounds each near relative error by
$2(n+1)e^{-(8/5)\sqrt n}$.
Using the global envelope and
$\sum_p\varepsilon^{d(p)}=2/(1-\varepsilon)$, the normalized error
between the actual sum \eqref{eq:bridge-exact} and the infinite ideal
sum is at most
\begin{equation}\label{eq:bridge-actual-error}
 \frac{456AB}{(1-\varepsilon)^2}
 \left((n+1)e^{-(8/5)\sqrt n}+\varepsilon^H\right)
 <\frac{AB}{1000H}.
\end{equation}
For the far part use the union of $d(p)\ge H$ and $d(q)\ge H$.
After multiplying by $H/AB$, the near bound decreases for
$\sqrt n\ge40$ and is below $10^7e^{-64}$ at the endpoint; the far
bound decreases for integer $H\ge7$. This proves the strict inequality.

For two successive mode events in different coordinates of one
rectangle, the exact event times are
$\max(\alpha,\beta)$ and $\min(\alpha,\beta)-\log\rho_H$.
If their gap is less than $4$, then
$|\alpha-\beta|+\log\rho_H<4$.
There is a unique zero $\widetilde\alpha$ of $F$ within $.001$ of
$\alpha$, and a unique zero $\widetilde\beta$ of $G$ within $.001$ of
$\beta$; both functions are increasing in the respective
$2.02$ neighborhoods. To verify this explicitly, put $\eta=e^{-12}$.
Coordinate concavity gives, with missing negative indices set to zero,
\[
 a_{2+j}\le A^2\eta(A\eta)^j,\qquad
 a_{-1-j}\le A^{-1}\eta(A^{-1}\eta)^j,\qquad j\ge0.
\]
If $|u-\alpha|\le T$ and $r_T=e^T\eta$, the two tail bounds are
\begin{align*}
 \sum_{p\notin\{0,1\}}a_p
 &\le\frac{\eta(e^T+e^{2T})}{1-r_T},\\
 \sum_{p\notin\{0,1\}}|p|a_p
 &\le\frac{\eta[e^T+e^{2T}(2-r_T)]}{(1-r_T)^2}.
\end{align*}
For $T=.001$, the first is below $.00002$, whereas $1-A$ has
opposite signs of magnitude exceeding $.000999$ at the endpoints.
For $T=2.02$, the second is below $.001$, so
$F'(u)>e^{-2.02}-.001>.128$. This proves the assertions directly.
The proof for $G$ is identical. These are exact Eulerian zeros:
the Frobenius identity expresses $F$ as the nonzero prefactor
$(-1)^K(HX)^K/K^n$ times
$A_n(-1/(HX))/(1+1/(HX))^{n+1}$.

Choose $u_*=(\widetilde\alpha+\widetilde\beta)/2$.
Its distances from the two balances are below $2.001$, and
$F(u_*)G(u_*)\le0$. Equations~\eqref{eq:bridge-product} and
\eqref{eq:bridge-actual-error} show that the normalized actual sum is
less than
\[
 -\frac{AB}{H}\left(\frac H{H+1}-.08-.001\right)<-.79AB/H<0.
\]
This sign is opposite to the two separators around the paired cluster.
The midpoint lies between those separators: the independent zeros
are within $.001$ of their balances, and the separators are at least
$2$ from the actual events.

\subsection{The phase at a separator}

\begin{lemma}\label{lem:bridge-phase}
Let $n>1600$, $u\ge3\sqrt n$, and let $(K,H)$ maximize $W_u$, with
$K,H\ge7$ and coordinate-neighbor ratios at most $e^{-2}$.
For $r=K+H-2$,
\[
 |\phi(-u)-(r-n)\pi|<3\pi/4.
\]
\end{lemma}

\begin{proof}
Set $C=1+e^u$, $z=-u+i\pi$, $N=n-K+1$, and
$H_n(y)=\sum_{\ell=0}^{n-1}A(n,\ell+1)y^\ell/(n)_\ell$.
Eulerian symmetry gives $R_{-u}(z)=n!C^nz^{-n-1}H_n(z/C)$.
By Lemma~\ref{lem:factorial}, $H_n$ has negative real roots, so its
coefficients tilted by a positive geometric sequence are log-concave.
Put $\varrho=|z|/C$ and
$c_\ell=A(n,\ell+1)\varrho^\ell/(n)_\ell$.

Comparison first in $k$ and then in $h$ gives $K,H<n/u+1$ and
\begin{equation}\label{eq:bridge-h-mode}
 N\log\rho_H\le u\le N\log\rho_{H-1}.
\end{equation}
The Eulerian relative error at $K$ is below $.001$, since
$n/K>3\sqrt n-9$ and $(n+1)e^{9-3\sqrt n}<.001$.
Writing $A_\pm$ for the $k$-neighbor ratios, we have
\begin{align*}
 \frac{c_K}{c_{K-1}}
 &\le\frac{A_+}{.999}\frac{H\sqrt{u^2+\pi^2}}N\frac XC,\\
 \frac{c_{K-2}}{c_{K-1}}
 &\le\frac{A_-}{.999}\frac{(N+1)C}{HX\sqrt{u^2+\pi^2}}.
\end{align*}
Here $u>120$, $N>1585$, and \eqref{eq:bridge-h-mode} gives
$Hu/N\le7/6$ and $(N+1)/(Hu)<1.001(8/7)$.
Also $C/X<1.001$ and $\sqrt{1+\pi^2/u^2}<1.001$.
Both ratios are therefore below $1.18e^{-2}<1/6$.
Log-concavity bounds the sum of the other terms by
$2\sum_{j\ge1}6^{-j}=2/5$ times the term of degree $K-1$.

That monomial dominates on $|y|=\varrho$. Along the upper semicircle
the remaining factor stays within distance $.4$ of $1$, so its
canonical argument has magnitude below $.42$. It starts positive,
which fixes the argument without a winding ambiguity. With
$\theta=\arctan(\pi/u)$ we obtain
\[
 \phi(-u)=-(N+1)(\pi-\theta)+\zeta,\qquad |\zeta|<.42.
\]
The logarithmic-mean bounds
$1/\log\rho_{H-1}\ge\sqrt{H(H-1)}>H-.6$ and
$1/\log\rho_H\le H+.5$ give $H-.6<N/u\le H+.5$.
Furthermore
\[
 0<(N+1)/u-(N+1)\theta/\pi
 <(n+1)\pi^2/(3u^3)<.004,\qquad 1/u<.009.
\]
Thus $H-.604<(N+1)\theta/\pi<H+.509$.
Subtracting $(r-n)\pi=-(N+1)\pi+H\pi$ bounds the phase error by
$.604\pi+.42<3\pi/4$.
\end{proof}

\subsection{Assembling the transition}

\begin{lemma}\label{lem:bridge}
For every $n>1600$ there are an integer $r<3\sqrt n/5$ and a number
$u\in[4\sqrt n-8,4\sqrt n]$ such that exactly $r$ zeros of $P_n$ lie
outside $|x|=e^u$, all are distinct and negative, and
\[
 \operatorname{sgn}P_n(-e^u)=(-1)^{n+r},\qquad
 |\phi(-u)-(r-n)\pi|<3\pi/4.
\]
\end{lemma}

\begin{proof}
Write $v=\sqrt n>40$ and $L=\lfloor\sqrt{n/(4\log n)}\rfloor$.
Lemma~\ref{lem:outer} supplies $2L$ distinct negative roots beyond
its innermost separator $X_L^-$.
Start at $u_{\rm out}=\log X_L^- -2$, with central cell $(L+1,L+1)$.
Its left ratios are below $e^{-2}/2$ and its right ratios below
$e^2/1000<e^{-2}$, by the outer lemma.
Since $\log n>25/4$, we have $7\le L\le v/5$ and $L+1\le9v/40$.
For $D=2(L+1+\rho_L^L)<2(L+4)$,
\[
 u_{\rm out}=n\log\rho_L-\log D-2
 >40v/9-\log(2v/5+8)-2>37v/9>4v.
\]
The logarithm comparison holds at $v=40$ and its right side has the
larger derivative thereafter. Lemma~\ref{lem:bridge-vertex} gives an
exact exterior count of $2L$, already supplied by the outer lemma.

We verify all geometry down to $u=4v-8$. A maximum of $W_u$ exists
on the infinite quadrant. Comparison first in $k$ gives $K<n/u+1<n$,
and then comparison in $h$ gives $H<n/u+1$.
At an event its stepping coordinate is less than $n/u$, directly
from \eqref{eq:bridge-balances}. Put $a=v^2/(4v-8)$.
Thus the stepping coordinate is less than $a$ and the other less
than $a+1$. The same-coordinate spacings exceed
\[
 \frac{v^2}{(a+1)^2},\qquad \frac{v^2-a}{(a+1)^2},
\]
both greater than $12$. For $v\ge41$, clearing denominators in the
second comparison gives
$4\{v^3(v-41)+18v^2+192v-192\}>0$.
For $40<v<41$, $a<11$, so the stepping coordinate is at most $10$
and the other at most $11$; the smaller spacing exceeds $1590/121>13$.
In a paired rectangle both stepping coordinates are less than $a$,
so its $\delta_k,\delta_h$ exceed $12$ as well.
Every mode center satisfies
\[
 \max(K,H)<a+1<v/4+2,\quad
 K+H<v/2+4<3v/5,\quad n-K-H>(K+H)^2.
\]
The mode path applies since its total degree remains below
$\lfloor\sqrt n\rfloor$. Its coordinates never decrease as $u$
decreases, and hence remain at least $L+1\ge8$.

Join consecutive events into a cluster when their gap is \emph{less}
than $4$; gaps at least $4$ separate clusters.
A cluster contains at most two events: three would put two
same-coordinate events within $8$. A pair uses both coordinates.
The spacing lower bounds depend only on the earlier event, so this
also excludes a hypothetical third event just below the eight-unit band.

Choose $u_{\rm in}\in[4v-8,4v]$ at distance at least $2$ from every
event. If $4v$ is already safe, use it. Otherwise take the cluster
with an event within $2$ of $4v$ and move to distance $2$ below its
lowest event. Its width is less than $4$, so the displacement is
less than $8$; the next cluster is at least $4$ away.
Use the two endpoints and the midpoints of intervening separating
gaps as separators. A coordinate neighbor of their maximizing cell
belongs to an adjacent total degree, so its ratio is at most $e^{-2}$.
Lemma~\ref{lem:bridge-vertex} applies at each separator.

Across a singleton cluster the total degree increases by one.
The opposite separator signs give one negative real zero.
Across a paired cluster it increases by two; the separator signs
agree and the intermediate sign established above gives two distinct
negative real zeros. The Rouch\'e counts give exactly one or two
zeros in the respective annuli. Starting with the $2L$ outer zeros,
all $r=K+H-2$ exterior zeros at the final separator are therefore
distinct and negative.
Finally $u_{\rm in}\ge4v-8>3v$, so Lemma~\ref{lem:bridge-phase}
gives the phase estimate, while the coordinate bound gives
$r<3v/5$.
\end{proof}

All decimal constants here are terminating rational numbers.
Their exponential comparisons and polynomial endpoint checks are also
certified by the accompanying standard-library verification scripts.
